\chapter[Literature Review]{Literature Review}
\label{Chap:Literature}

% ***************************************************
% Literature Review
% ***************************************************

% ********* Enter your text below this line: ********

\section{LiDAR SLAM Foundations}
Simultaneous localisation and mapping (SLAM) methods integrate perception and state estimation to produce globally consistent trajectories for mobile robots operating in uncertain environments \cite{cadena_2016_past}. LiDAR-centric pipelines exploit dense range observations that are resilient to illumination changes and long-range drift, but they must reconcile point-cloud registration, inertial fusion, and loop-closure hypotheses in real time. LIO-SAM exemplifies this design by combining deskewed scan registration with factor-graph smoothing and IMU preintegration in a tightly coupled estimator \cite{shan_2020_liosam}. Scan Context descriptors further extend the architecture, enabling loop closures through polar signatures that match previously visited scenes without resorting to heavy point-cloud alignment \cite{gkim-2018-iros}.

SC-LIO-SAM maintains these components while exposing a broad set of hyperparameters that govern everything from range-image resolution to downsampling voxel sizes and feature-count thresholds. Because each module contributes to both computational load and Jacobian conditioning, parameter selection directly impacts accuracy and latency. Field deployments therefore hinge on well-chosen settings that keep the surf and corner feature sets diverse enough for stable scan-to-map optimisation without overwhelming the solver.

\section{Hyperparameter Tuning Strategies}
Historically, SLAM practitioners tuned voxel filters, feature minima, and gating thresholds by trial-and-error using domain intuition, but this approach is brittle when geometries or sensor rigs change. Automated strategies have emerged to search these spaces systematically. Koide et al. cast LiDAR odometry tuning as a black-box optimisation problem and leveraged Gaussian-process surrogates to recommend parameters that improved pose accuracy without inspecting internal gradients \cite{koide_2021_adaptive}. Goudreault et al. formulated LiDAR-in-the-loop hyperparameter optimisation, repeatedly replaying datasets to evaluate candidate parameters and ranking them via performance metrics gathered from the full estimator stack \cite{goudreault_2023_lidarintheloop}. These techniques demonstrate that LiDAR odometry quality can be meaningfully improved by targeted exploration of voxel-grid and noise-related parameters, yet they typically operate offline and do not react to intra-trajectory variability.

The present project adopted similar principles while executing one-factor-at-a-time sweeps on SC-LIO-SAM. By replaying identical MulRan runs under different `mapping\_surf\_leaf\_size` settings, we confirmed that surf voxel size dominated the trade-off between scan alignment fidelity and runtime cost. This evidence positions `mapping\_surf\_leaf\_size` as the most leverageable online control input and motivates a controller that can adapt it mid-trajectory rather than selecting a single static value.

\section{Reinforcement Learning for Perception Control}
Reinforcement learning (RL) methods can observe environment feedback, generate temporally correlated actions, and therefore adapt perception systems during operation. Messikommer et al. framed visual odometry tuning as an RL problem where an attention-equipped PPO agent modulated exposure-like parameters on the fly using latent token observations \cite{nicomessikommer_2024_reinforcement}. Complementary work has adjusted VO camera settings and visual SLAM thresholds through deep RL policies that react to perceived brightness or feature statistics \cite{zhang_2024_efficient, li_2025_uampc}. Beyond cameras, researchers have begun to condition LiDAR perception on RL-controlled parameters: Kabirat Bolanle Olayemi et al. tuned LiDAR configurations for navigation policies, and Ma et al. developed DQN variants that adapt SLAM parameters based on historical observations \cite{kabiratbolanleolayemi_2023_the, ma_2025_mambadqn}. Collectively, these studies highlight RL’s ability to convert high-dimensional sensor diagnostics into actionable control signals for perception systems.

Despite this progress, existing RL-perception controllers seldom integrate with full LiDAR SLAM stacks that rely on ROS-based orchestration, ground-truth evaluation tools, and runtime parameter servers. Most approaches either act on simplified simulators or treat the SLAM pipeline as a static black box, which limits their ability to exploit fine-grained metrics such as feature counts, odometry cadence, or action history. A LiDAR-specific RL environment must therefore ingest system-level metrics, maintain safe parameter bounds, and coordinate with replay infrastructure like MulRan \cite{kim_2020_mulran}.

\section{Gap Analysis}
The reviewed literature establishes two trends: LiDAR SLAM quality is highly sensitive to voxel-filter and feature-selection parameters, and RL policies can adapt perception modules when fed informative observations. However, no existing work evaluates whether an RL agent can continuously tune SC-LIO-SAM’s `mapping\_surf\_leaf\_size` while replaying real LiDAR sequences, nor do prior efforts couple PPO training loops to ROS-native stacks that emit both odometry trajectories and auxiliary metrics. Bridging this gap requires an environment that grounds RL decisions in real scalar measurements produced by SC-LIO-SAM, leverages PPO to learn policies from noisy MulRan replays, and measures progress through evo-derived pose errors. This thesis therefore targets the underexplored intersection between LiDAR SLAM hyperparameter sensitivity and adaptive RL control, delivering evidence for or against RL’s ability to stabilise SC-LIO-SAM across diverse MulRan traversals \cite{shan_2020_liosam,kim_2020_mulran,nicomessikommer_2024_reinforcement}.
